% !TEX root = ../main.tex
\chapter{HYPOTHESIS}\label{hypothesis}

\textbf{H1 — Mechanical Design:}\label{H1}
The PRANK parallel robotic structure, as defined by its SolidWorks CAD model and kinematic analysis, will achieve a functional workspace covering at least $\pm25^{\circ}$ of dorsiflexion/plantarflexion, $\pm15^{\circ}$ of inversion/eversion, and $\pm10^{\circ}$ of internal/external rotation — values established as the minimum therapeutic range of motion required for ankle rehabilitation \cite{Brockett2016, Tsoi2009}.

\textbf{H2 — Control System:}\label{H2}
The admittance-based control system implemented in the computational simulation environment will demonstrate stable and responsive behaviour across passive, assistive, and resistive rehabilitation modes, with a settling time under 500\,ms and a force tracking error below 15\% of the commanded torque, ensuring safe and adaptive interaction with the patient throughout therapy \cite{Zhang2018, Li2020passive}.

\textbf{H3 — Hardware-in-the-Loop Validation:}\label{H3}
The hardware-in-the-loop configuration — in which the STM32 development board executes the admittance control firmware in real time while exchanging state and command data with the computational simulator — will reproduce the control trajectories obtained in the pure software simulation with a root mean square deviation below $2^{\circ}$ and a communication latency compatible with real-time operation ($<10$\,ms cycle), demonstrating the feasibility of the embedded implementation before physical fabrication.
